\section{Introduction}
\label{sec:intro}

Much applied computer-vision work starts from a public dataset that someone has already augmented: to reach a target size, every photograph ships with rotated, flipped, recoloured and filtered copies, and nothing records which copies belong together (Figure~\ref{fig:concept}a). The user then draws a random 80/20 split, trains, and reports test accuracy.

This paper shows that the number obtained this way measures very little, and that it can be wrong in either direction. The first error is leakage. With $k$ copies per photograph and test fraction $\rho$, a test image has a sibling in the training set with probability $\approx1-\rho^{k-1}$ (Proposition~\ref{prop:exposure}). At $k=4$ that is $99.2\%$ of the test set; at $k=16$ it is all of it, and the test score becomes, up to a weight of $3\times10^{-11}$, an accuracy on images the model has already seen in another guise (Proposition~\ref{prop:decomp}). The second error is that the test images are themselves edited copies, so even a leak-free split scores the model on filtered images that no deployment will see. On a controlled release built from CIFAR-10, a frozen-feature 1-NN classifier reports \numWeakKnnRep{} accuracy where the same model reaches \numWeakKnnFresh{} on fresh photographs; with stronger edits, the reported number falls \emph{below} the true one (Figure~\ref{fig:inflation}).

The obvious repair, recovering the groups with a near-duplicate detector, is harder than it looks. Released copies are deliberately dissimilar: an edge map of a leaf is closer, in an ImageNet embedding, to the edge maps of \emph{other} leaves than to its own original. And removing leakage needs complete groups: a group cut into two fragments still leaks with probability $1-\rho$ (Proposition~\ref{prop:frag}). On two real releases with known groups, splitting by the groups found by any detector we tried, including ours, leaves more than $90\%$ of test images with a sibling in training.

Yet the damage is not where nominal exposure suggests. Inflation is concentrated on test images whose closest training sibling is close in the model's own representation; siblings that are hard to match are also hard to exploit (Figure~\ref{fig:exploit}). A detector therefore does not need complete groups to be useful: splitting by the groups it finds removes most of the inflation when its similarity views match the release's edits (Section~\ref{sec:protocol}).

\paragraph{Contributions.}
\begin{itemize}[leftmargin=1.2em,itemsep=1pt,topsep=2pt]
\item \textbf{Theory of pre-augmented splits.} Exact exposure probabilities, the decomposition of reported accuracy into seen and unseen parts, and a fragmentation result showing why pairwise de-duplication does not remove leakage (Section~\ref{sec:theory}).
\item \textbf{Measurement.} On a controlled release with fresh ground-truth test images and on two real releases, we measure inflation across augmentation strength, copies per image, data size, classifier and backbone, and show the error can have either sign (Section~\ref{sec:inflation}).
\item \textbf{AugLeak}, a metadata-free sibling detector with an approximately calibrated error rate: a structure view robust to photometric filters, transform search, hubness correction and an extreme-value null derived from class labels (Section~\ref{sec:method}). On both real releases with ground-truth groups it has the highest average precision of all methods we tested (Section~\ref{sec:detect}).
\item \textbf{An evaluation of metadata-free protocols} showing when detected-group splits recover the oracle accuracy and why test filtering does not, and an \textbf{audit} of six public collections that finds repeat shots filed as independent photographs, test images duplicated in training (some under another label; half of one MRI test set) and radiographs filed under two patient identifiers (Sections~\ref{sec:protocol}--\ref{sec:audit}). Code and an auditing tool are released.
\end{itemize}
